\documentclass[10pt,a4paper]{article}

\usepackage[margin=1.05in]{geometry}
\usepackage{setspace}
\usepackage{amsmath,amssymb,amsthm}
\usepackage{graphicx}
\usepackage{booktabs}
\usepackage{array}
\usepackage{caption}
\usepackage{enumitem}
\usepackage{titlesec}
\usepackage[expansion=false]{microtype}
\usepackage{seqsplit}
\usepackage{hyperref}
\usepackage{fancyhdr}

\hypersetup{colorlinks=true,linkcolor=black,urlcolor=black,citecolor=black}
\captionsetup{font=small,labelfont=bf,skip=4pt,justification=raggedright}

\titleformat{\section}{\normalsize\bfseries}{}{0em}{}
\titleformat{\subsection}{\small\bfseries}{}{0em}{}
\titlespacing*{\section}{0pt}{1.35em}{0.5em}
\titlespacing*{\subsection}{0pt}{0.95em}{0.35em}

\setstretch{1.14}
\setlist[itemize]{leftmargin=1.15em,itemsep=0.1em,topsep=0.2em}
\setlength{\parindent}{1.2em}

\pagestyle{fancy}
\fancyhf{}
\fancyhead[L]{\small\itshape Elements of Position}
\fancyhead[R]{\small\itshape Companion: Verification and Models}
\fancyfoot[C]{\small\thepage}
\renewcommand{\headrulewidth}{0pt}

\newcommand{\code}[1]{\texttt{#1}}

\begin{document}

\begin{center}
\vspace*{1.0em}
{\large Elements of Position}\\[1.0em]
{\LARGE\bfseries Companion: Verification and Models}\\[1.4em]
{\normalsize Justin Erholtz}\\[0.3em]
{\small 13 September 2026}
\end{center}

\vspace{1.2em}

\noindent
\textit{A note.}
This is not a numbered part of the series. It is documentation: a
description of the code that verifies claims made in Papers I--III at a
depth their own printed tables cannot show, and that gives concrete visual
form to the plan/section/elevation reading Paper IV names but does not
itself draw. Nothing here is a philosophical argument, and nothing here
extends the series' claims. Where this companion's tooling turned up
structure unrelated to anything the papers argue, that material is
deliberately not included here; it is held in a separate, clearly labeled
exploratory record.

\vspace{0.8em}

\section{Why this exists}

Papers I through III make claims that hold at every finished $n$: the
geometric mean locks at $1$, the marriage of Process A and Process B is
exact only at $n=5$, the Fibonacci docking error shrinks as the count
runs. Each paper demonstrates this with a printed table running to
$n=16$ --- a table can't run further and still be read. That is a real
limitation, not a rhetorical one: ``one run is every run'' (Paper III)
is a claim about all $n$, and sixteen rows support it by example, not
by exhaustion.

This companion exists to close that gap: two scripts that carry every
claim out to a depth no page could hold, self-test the results on every
execution, and give a live, orbitable model of the two-process geometry
Paper III describes only in prose and static figures.

\section{What is included, and what is deliberately not}

Two scripts, described in \S3--\S4. Both are built on formulas ported
verbatim from \texttt{\seqsplit{Elements\_of\_Position\_-\_Drawings.py}},
the script every paper already cites --- this companion is a second
reading of the same geometry, not a second geometry.

During the construction of these tools, several structures were found
that are true, checkable, and unrelated to anything the papers claim ---
among them a connection between the vertex structure explored in Model~2
and Euler's totient function, and a further connection to cyclotomic
polynomials. Those findings are real and documented, but held in a
separate exploratory record: they support nothing written in Papers I--V.
Keeping that boundary explicit is itself part of the discipline the
papers ask of themselves.

\section{Verification.py --- Papers I, II, III}

Ports \code{row\_for\_n}, \code{count\_process}, and \code{phi\_process}
directly from the drawings script, then runs the master sequence to
$n=10{,}000$ by default (configurable), producing:

\begin{itemize}
\item \textbf{\code{measurements.csv}} --- one row per $n$: product,
geometric mean, arithmetic mean, cross-ratio $-n$, half-walk $n\pi/2$,
chord, apothem, and \code{marriage\_gap} $= 2\cos(\pi/n) - \varphi$
(Paper III's closing lemma, generalized across the full range).
\item \textbf{\code{docking\_events.csv}} --- the Fibonacci docking table
from Paper III, Proposition 1, computed two ways per row: the paper's own
formula ($\tfrac12|F_k\Phi - F_{k-1}|$) and a closed form derived here via
Binet's identity, $\Phi^{k+2}/2$, confirmed to agree to 44+ decimal
digits under arbitrary-precision arithmetic. The paper's formula
subtracts two large, nearly equal numbers --- ordinary floating-point
cancellation, visible directly in the \code{cancellation\_error} column,
growing from $\sim 10^{-17}$ at $k=1$ to $\sim 10^{-13}$ by $k=18$. The
closed form has no such term and is used as the reported value.
\item \textbf{\code{live\_model.html}} --- Process A and Process B in one
orbitable 3D scene, with a configurable arrival point for B, Fibonacci
dockings marked, and the $n=5$ marriage called out directly.
\end{itemize}

\begin{table}[h]
\centering
\small
\begin{tabular}{@{}rrr@{}}
\toprule
$k$ & $F_k$ & docking error \\
\midrule
1 & 2 & 0.118034 \\
5 & 13 & 0.017221 \\
10 & 144 & 0.001553 \\
15 & 1597 & 0.000140 \\
18 & 6765 & 0.000033 \\
\bottomrule
\end{tabular}
\caption{A sample of the docking table at $n=10{,}000$ depth. Full table
in \code{docking\_events.csv}.}
\end{table}

Self-test on every run: \code{docking\_error == Phi**(k+2)/2} across all
rows in range, to within $10^{-12}$.

\section{Local Circle Models.py --- Papers I, III, IV}

Three models, matching Paper IV's own division of one object into plan,
section, and elevation:

\begin{description}[leftmargin=1.3em]
\item[Model 1 --- local circle.] One $n$-gon at a time, built raw from
the master sequence: radius $n/2$, chord $n\sin(\pi/n)$, apothem
$(n/2)\cos(\pi/n)$, central angle $2\pi/n$, and the isosceles spoke
triangle's base angles. This is Paper III's worked $n=5$ example (radius
$2.5$, central angle $72^\circ$, base angles $54^\circ$ --- reproduced
exactly by this script) generalized to any finished $n$.
\item[Model 2 --- plan overlay.] Every local circle from Model 1,
overlaid concentrically at one center --- the flat nesting Paper III
describes when it says scale does not matter.
\item[Model 3 --- elevation stack.] The same $n$-gons, stacked by height
$n/2$ instead of flattened to one plane. This height is not an arbitrary
choice: Paper I's own cone construction already lifts the own-origin to
height $n/2$, forming the $45$--$45$--$90$ triangle with vertices
$(0,0)$, $(n,0)$, $(n/2,n/2)$. The vertex line running straight down the
elevation stack at $\theta=0$ \emph{is} that triangle's hypotenuse,
already proved in Paper I. The section cut at the mid-edge angle traces
the apothem, $(n/2)\cos(\pi/n)$, which converges onto that same line as
$n$ grows --- a picture of the Archimedes inequality
$n\sin(\pi/n) < \pi < n\tan(\pi/n)$ that Paper I cites in its literature
review but does not itself draw.
\end{description}

Self-tests on every run: the relative gap $1 - (\text{apothem}/\text{radius})$
matches $\pi^2/(2n^2)$ to within 1\% for $n \geq 500$; the inscribed and
circumscribed Archimedes bounds hold on every side, for every $n$, with
no exceptions.

\section{Reproducing this companion}

\begin{verbatim}
python "Elements of Position - Verification.py" --nmax 10000
python "Elements of Position - Local Circle Models.py" --nmax 1000
\end{verbatim}

Both scripts print a \texttt{self-test:} block on completion. A run that
does not print all \texttt{PASS} lines should not be trusted, and neither
should any figure produced by it.

\section{Scope}

This companion verifies. It does not extend. Nothing here should be read
as a sixth part of the series, and nothing here bears on Paper IV's
analytic material (the derived-unit, vacuum/trace, and midline
discussion) --- that content is explicitly a naming choice within the
paper itself, not a numerical claim, and no tool in this companion
touches it either way.

\end{document}
